
\chapter{2017年江苏卷}

\section{填空题}

\begin{problem}
已知集合 \(A=\{1,2\}\)，\(B=\{a,a^2+3\}\). 若 \(A\cap B=\{1\}\)，则实数 \(a\) 的值为\fillinblank{}.
\end{problem}

\begin{answer}
\(1\)
\end{answer}

\begin{solution}
因为 \(A\cap B=\{1\}\)，所以 \(1\in B\). 因而 \(a=1\) 或 \(a^2+3=1\). 后一种情况给出 \(a^2=-2\)，不可能；所以 \(a=1\). 此时 \(B=\{1,4\}\)，确有 \(A\cap B=\{1\}\)，故填 \(1\).
\end{solution}



\begin{problem}
已知复数 \(z=(1+\i)(1+2\i)\)，其中 \(\i\) 是虚数单位，则 \(z\) 的模是\fillinblank{}.
\end{problem}

\begin{answer}
\(\sqrt{10}\)
\end{answer}

\begin{solution}
由复数乘法得 \(z=(1+\i)(1+2\i)=1+3\i+2\i^2=-1+3\i\). 因此
\[
|z|=\sqrt{(-1)^2+3^2}=\sqrt{10}.
\]
\end{solution}



\begin{problem}
某工厂生产甲，乙，丙，丁四种不同型号的产品，产量分别为 \(200\)，\(400\)，\(300\)，\(100\) 件. 为检验产品的质量，现用分层抽样的方法从以上所有的产品中抽取 \(60\) 件进行检验，则应从丙种型号的产品中抽取\fillinblank{} 件.
\end{problem}

\begin{answer}
\(18\)
\end{answer}

\begin{solution}
四种型号产品的总产量为 \(200+400+300+100=1000\) 件，分层抽样的抽样比为 \(60/1000\). 因而应从丙种型号产品中抽取
\[
300\times\frac{60}{1000}=18
\]
件.
\end{solution}



\begin{problem}
如图是一个算法流程图：若输入 \(x\) 的值为 \(\dfrac{1}{16}\)，则输出 \(y\) 的值是\fillinblank{}.

\bitmapfigure[max width=6.0cm]{img/2017/jiangsu/q4_fig1.png}
\end{problem}

\begin{answer}
\(-2\)
\end{answer}

\begin{solution}
因为输入值 \(x=\dfrac1{16}\le1\)，流程进入分支 \(y=2+\log_2x\). 所以
\[
y=2+\log_2\frac1{16}=2-4=-2.
\]
\end{solution}



\begin{problem}
若 \(\tan\!\left(\alpha-\dfrac{\pi}{4}\right)=\dfrac16\)，则 \(\tan\alpha=\fillinblank{}\).
\end{problem}

\begin{answer}
\(\dfrac75\)
\end{answer}

\begin{solution}
设 \(t=\tan\alpha\). 由正切的差角公式，得
\[
\frac{t-1}{1+t}=\frac16.
\]
解得 \(6t-6=1+t\)，即 \(5t=7\)，所以 \(\tan\alpha=t=\dfrac75\).
\end{solution}



\begin{problem}
如图，在圆柱 \(O_1O_2\) 内有一个球 \(O\)，该球与圆柱的上，下底面及母线均相切，记圆柱 \(O_1O_2\) 的体积为 \(V_1\)，球 \(O\) 的体积为 \(V_2\)，则 \(\dfrac{V_1}{V_2}\) 的值是\fillinblank{}.

\bitmapfigure[max width=6.0cm]{img/2017/jiangsu/q6_fig1.png}
\end{problem}

\begin{answer}
\(\dfrac32\)
\end{answer}

\begin{solution}
设球的半径为 \(r\). 球与圆柱的上、下底面相切，所以圆柱的高为 \(2r\)；球与圆柱的母线相切，所以圆柱的底面半径也为 \(r\). 因而
\[
V_1=\pi r^2\cdot2r=2\pi r^3,\qquad V_2=\frac43\pi r^3.
\]
所以
\[
\frac{V_1}{V_2}=\frac{2\pi r^3}{\frac43\pi r^3}=\frac32.
\]
\end{solution}



\begin{problem}
记函数 \(f(x)=\sqrt{6+x-x^2}\) 的定义域为 \(D\). 在区间 \([-4,5]\) 上随机取一个数 \(x\)，则 \(x\in D\) 的概率是\fillinblank{}.
\end{problem}

\begin{answer}
\(\dfrac59\)
\end{answer}

\begin{solution}
要使根式有意义，必须有 \(6+x-x^2\ge0\)，即
\[
(x+2)(x-3)\le0.
\]
因此 \(D=[-2,3]\). 在 \([-4,5]\) 上等可能取数时，事件 \(x\in D\) 对应区间长度为 \(5\)，样本区间长度为 \(9\)，故所求概率为
\[
\frac{3-(-2)}{5-(-4)}=\frac59.
\]
\end{solution}



\begin{problem}
在平面直角坐标系 \(xOy\) 中，双曲线 \(\dfrac{x^2}{3}-y^2=1\) 的右准线与它的两条渐近线分别交于点 \(P\)，\(Q\)，其焦点是 \(F_1\)，\(F_2\)，则四边形 \(F_1PF_2Q\) 的面积是\fillinblank{}.
\end{problem}

\begin{answer}
\(2\sqrt{3}\)
\end{answer}

\begin{solution}
双曲线的标准参数为 \(a^2=3\)，\(b^2=1\)，所以 \(c=\sqrt{a^2+b^2}=2\). 右准线方程为 \(x=\dfrac{a^2}{c}=\dfrac32\)，两条渐近线为 \(y=\pm\dfrac{x}{\sqrt3}\). 因此
\[
P=\left(\frac32,\frac{\sqrt3}{2}\right),\qquad Q=\left(\frac32,-\frac{\sqrt3}{2}\right).
\]
双曲线的两个焦点为 \(F_1=(-2,0)\)，\(F_2=(2,0)\). 四边形 \(F_1PF_2Q\) 的两条对角线 \(F_1F_2\) 与 \(PQ\) 互相垂直，且长度分别为 \(4\) 与 \(\sqrt3\)，所以面积为
\[
\frac12\cdot4\cdot\sqrt3=2\sqrt3.
\]
\end{solution}



\begin{problem}
等比数列 \(\{a_n\}\) 的各项均为实数，其前 \(n\) 项和为 \(S_n\)，已知 \(S_3=\dfrac74\)，\(S_6=\dfrac{63}{4}\)，则 \(a_8=\fillinblank{}\).
\end{problem}

\begin{answer}
\(32\)
\end{answer}

\begin{solution}
设等比数列的公比为 \(q\). 前三项与后三项分别相差三个位置，所以
\[
S_6=S_3+q^3S_3=S_3(1+q^3).
\]
由已知得 \(1+q^3=\dfrac{S_6}{S_3}=9\)，从而 \(q=2\). 又 \(S_3=a_1(1+q+q^2)=7a_1=\dfrac74\)，故 \(a_1=\dfrac14\). 因此
\[
a_8=a_1q^7=\frac14\cdot2^7=32.
\]
\end{solution}



\begin{problem}
某公司一年购买某种货物 \(600\) 吨，每次购买 \(x\) 吨，运费为 \(6\) 万元/次，一年的总存储费用为 \(4x\) 万元. 要使一年的总运费与总存储费用之和最小，则 \(x\) 的值是\fillinblank{}.
\end{problem}

\begin{answer}
\(30\)
\end{answer}

\begin{solution}
因为每次购买 \(x\) 吨，一年购买 \(600\) 吨，所以购买次数为 \(600/x\)，其中 \(x>0\). 一年的总费用为
\[
T(x)=6\cdot\frac{600}{x}+4x=\frac{3600}{x}+4x.
\]
由基本不等式，
\[
T(x)=\frac{3600}{x}+4x\ge2\sqrt{\frac{3600}{x}\cdot4x}=240.
\]
当且仅当 \(\dfrac{3600}{x}=4x\) 时取等号，即 \(x^2=900\). 因为 \(x>0\)，所以 \(x=30\).
\end{solution}



\begin{problem}
已知函数 \(f(x)=x^3-2x+\e^x-\dfrac{1}{\e^x}\)，其中 \(\e\) 是自然对数的底数. 若 \(f(a-1)+f(2a^2)\le 0\)，则实数 \(a\) 的取值范围是\fillinblank{}.
\end{problem}

\begin{answer}
\(\left[-1,\dfrac12\right]\)
\end{answer}

\begin{solution}
函数 \(f\) 是奇函数，因为 \(x^3-2x\) 与 \(\e^x-\e^{-x}\) 都是奇函数. 且
\[
f'(x)=3x^2-2+\e^x+\e^{-x}=3x^2+\left(\e^{x/2}-\e^{-x/2}\right)^2\ge0,
\]
等号只在 \(x=0\) 时成立，所以 \(f\) 在 \(\R\) 上严格递增.

令 \(u=a-1\)，则由奇性有 \(-f(u)=f(-u)=f(1-a)\). 因此
\[
f(a-1)+f(2a^2)\le0
\Longleftrightarrow f(2a^2)\le f(1-a).
\]
由 \(f\) 严格递增，以上不等式等价于 \(2a^2\le1-a\)，即
\[
2a^2+a-1\le0.
\]
解得 \(a\in\left[-1,\dfrac12\right]\).
\end{solution}



\begin{problem}
如图，在同一个平面内，向量 \(\overrightarrow{OA}\)，\(\overrightarrow{OB}\)，\(\overrightarrow{OC}\) 的模分别为 \(1\)，\(1\)，\(\sqrt2\)，\(\overrightarrow{OA}\) 与 \(\overrightarrow{OC}\) 的夹角为 \(\alpha\)，且 \(\tan\alpha=7\)，\(\overrightarrow{OB}\) 与 \(\overrightarrow{OC}\) 的夹角为 \(45^\circ\). 若 \(\overrightarrow{OC}=m\overrightarrow{OA}+n\overrightarrow{OB}\)（\(m\)，\(n\in\R\)），则 \(m+n=\fillinblank{}\).

\bitmapfigure[max width=6.0cm]{img/2017/jiangsu/q12_fig1.png}
\end{problem}

\begin{answer}
\(3\)
\end{answer}

\begin{solution}
由 \(\tan\alpha=7\) 且 \(\alpha\) 为图中两向量的夹角，得
\[
\cos\alpha=\frac1{\sqrt{50}}=\frac1{5\sqrt2},\qquad \sin\alpha=\frac7{\sqrt{50}}=\frac7{5\sqrt2}.
\]
图中 \(\overrightarrow{OA}\) 与 \(\overrightarrow{OB}\) 分居 \(\overrightarrow{OC}\) 两侧，且 \(\angle BOC=45^\circ\)，所以
\[
\overrightarrow{OA}\cdot\overrightarrow{OB}=\cos(\alpha+45^\circ)
=\frac{1-7}{10}=-\frac35.
\]
又因为 \(|\overrightarrow{OC}|=\sqrt2\)，故
\[
\overrightarrow{OC}\cdot\overrightarrow{OA}=\sqrt2\cos\alpha=\frac15,qquad
\overrightarrow{OC}\cdot\overrightarrow{OB}=\sqrt2\cos45^\circ=1.
\]
将 \(\overrightarrow{OC}=m\overrightarrow{OA}+n\overrightarrow{OB}\) 分别与 \(\overrightarrow{OA}\)、\(\overrightarrow{OB}\) 作数量积，得
\[
m-\frac35n=\frac15,qquad -\frac35m+n=1.
\]
解得 \(m=\dfrac54\)，\(n=\dfrac74\)，所以 \(m+n=3\).
\end{solution}



\begin{problem}
在平面直角坐标系 \(xOy\) 中，\(A(-12,0)\)，\(B(0,6)\)，点 \(P\) 在圆 \(O:x^2+y^2=50\) 上. 若 \(\overrightarrow{PA}\cdot \overrightarrow{PB}\le 20\)，则点 \(P\) 的横坐标的取值范围是\fillinblank{}.
\end{problem}

\begin{answer}
\(\left[-5\sqrt{2},1\right]\)
\end{answer}

\begin{solution}
设 \(P=(x,y)\). 则
\[
\overrightarrow{PA}\cdot\overrightarrow{PB}=(-12-x,-y)\cdot(-x,6-y)=x^2+y^2+12x-6y=50+12x-6y.
\]
条件等价于 \(12x-6y\le-30\)，即 \(y\ge2x+5\).

由 \(P\) 在圆 \(x^2+y^2=50\) 上，必有 \(x\ge-\sqrt{50}=-5\sqrt2\). 若 \(x>1\)，则 \(2x+5>0\)，由 \(y\ge2x+5\) 可平方得
\[
50-x^2=y^2\ge(2x+5)^2,
\]
从而 \(5x^2+20x-25\le0\)，即 \((x-1)(x+5)\le0\)，这与 \(x>1\) 矛盾，所以 \(x\le1\).

反之，任取 \(x\in[-5\sqrt2,1]\)，取圆的上半圆点 \(y=\sqrt{50-x^2}\). 当 \(x\le-\dfrac52\) 时，\(2x+5\le0\le y\)；当 \(-\dfrac52<x\le1\) 时，\(2x+5>0\)，且
\[
50-x^2-(2x+5)^2=5(x+5)(1-x)\ge0,
\]
故 \(y\ge2x+5\). 因而所有 \(x\in[-5\sqrt2,1]\) 都能取到，所求范围为
\[
[-5\sqrt2,1].
\]
\end{solution}



\begin{problem}
设 \(f(x)\) 是定义在 \(\R\) 上且周期为 \(1\) 的函数，在区间 \([0,1)\) 上，\(f(x)=\begin{cases}
x^2,&x\in D,\\
x,&x\notin D,
\end{cases}\) 其中 \(D=\left\{x\mid x=\dfrac{n-1}{n},\ n\in \N^*\right\}\)，则方程 \(f(x)-\lg x=0\) 的解的个数是\fillinblank{}.
\end{problem}

\begin{answer}
\(8\)
\end{answer}

\begin{solution}
设 \(x=k+t\)，其中 \(k=\lfloor x\rfloor\) 为非负整数，\(0\le t<1\). 当 \(0<x<1\) 时，\(\lg x<0\)，而 \(f(x)=f(t)\ge0\)，所以没有解. 当 \(x\ge10\) 时，\(\lg x\ge1\)，而 \(0\le f(t)<1\)，也没有解. 因此只需考虑 \(k=1,2,\ldots,9\).

先看 \(t=0\). 此时 \(f(t)=0\)，方程要求 \(\lg k=0\)，故 \(k=1\)，得到解 \(x=1\).

若 \(0<t<1\) 且 \(t\in D\)，则 \(t=1-\dfrac1n\)（\(n\ge2\)），所以 \(t\) 和 \(t^2\) 都是介于 \(0\) 与 \(1\) 之间的有理数. 无论 \(f(t)=t\) 还是 \(f(t)=t^2\)，若方程成立，都有 \(k+t=10^{f(t)}\). 但对任意 \(0<r<1\) 的有理数 \(r\)，\(10^r\) 是无理数，这与 \(k+t\) 为有理数矛盾. 因此所有 \(0<t<1\) 的解都满足 \(t\notin D\)，从而 \(f(t)=t\).

对每个 \(k\ge1\)，令 \(h_k(t)=\lg(k+t)-t\). 因为
\[
h_k'(t)=\frac1{(k+t)\ln10}-1<0,
\]
所以 \(h_k\) 严格递减. 当 \(k=1\) 时，\(h_1(0)=0\)，除 \(t=0\) 外没有根；当 \(k=2,3,\ldots,8\) 时，\(h_k(0)=\lg k>0\)，且 \(h_k(1)=\lg(k+1)-1<0\)，所以每个 \(k\) 恰有一个 \(t\in(0,1)\) 使 \(h_k(t)=0\)，共得到 \(7\) 个解. 当 \(k=9\) 时，\(h_9(0)>0\)，\(h_9(1)=\lg10-1=0\)，故在 \(0\le t<1\) 内没有根.

综上，方程共有 \(1+7=8\) 个解.
\end{solution}


\section{解答题}

\begin{problem}
如图，在三棱锥 \(A-BCD\) 中，\(AB\perp AD\)，\(BC\perp BD\)，平面 \(ABD\perp\) 平面 \(BCD\)，点 \(E\)，\(F\)（\(E\) 与 \(A\)，\(D\) 不重合）分别在棱 \(AD\)，\(BD\) 上，且 \(EF\perp AD\).

求证：

\begin{enumerate}
\item \(EF\parallel\) 平面 \(ABC\)；
\item \(AD\perp AC\).
\end{enumerate}

\bitmapinlinefigure[8.1cm][]{img/2017/jiangsu/q15_fig1.png}
\end{problem}

\begin{solution}
\begin{enumerate}
\item 直线 \(EF\) 在平面 \(ABD\) 内，且 \(EF\perp AD\). 又因为 \(AB\perp AD\)，所以在平面 \(ABD\) 内，过点 \(E\) 的直线 \(EF\) 与 \(AB\) 平行，即 \(EF\parallel AB\). 由于 \(AB\) 在平面 \(ABC\) 内，且 \(EF\) 不在平面 \(ABC\) 内，因此 \(EF\parallel\) 平面 \(ABC\).
\item 平面 \(ABD\perp\) 平面 \(BCD\)，两平面的交线为 \(BD\). 又 \(BC\subset\) 平面 \(BCD\)，且 \(BC\perp BD\)，所以 \(BC\perp\) 平面 \(ABD\)，从而 \(BC\perp AD\). 已知 \(AB\perp AD\)，而 \(AB\)，\(BC\) 是平面 \(ABC\) 内相交的两条直线，所以 \(AD\perp\) 平面 \(ABC\)，故 \(AD\perp AC\).
\end{enumerate}
\end{solution}


\begin{problem}
已知向量 \(\bs{a}=(\cos x,\sin x)\)，\(\bs{b}=(3,-\sqrt3)\)，\(x\in[0,\pi]\).

\begin{enumerate}
\item 若 \(\bs{a}\parallel \bs{b}\)，求 \(x\) 的值；
\item 记 \(f(x)=\bs{a}\cdot \bs{b}\)，求 \(f(x)\) 的最大值和最小值以及对应的 \(x\) 的值.
\end{enumerate}
\end{problem}

\begin{solution}
\begin{enumerate}
\item 由 \(\bs{a}\parallel\bs{b}\)，得
\[
\begin{vmatrix}\cos x&\sin x\\3&-\sqrt3\end{vmatrix}=0,
\]
即 \(\sqrt3\cos x+3\sin x=0\). 因为 \(x\in[0,\pi]\)，所以 \(x\) 在第二象限，且 \(\tan x=-\dfrac{1}{\sqrt3}\)，故 \(x=\dfrac{5\pi}{6}\).
\item \(f(x)=\bs{a}\cdot\bs{b}=3\cos x-\sqrt3\sin x=2\sqrt3\cos\left(x+\dfrac{\pi}{6}\right)\). 当 \(x\in[0,\pi]\) 时，\(x+\dfrac{\pi}{6}\in\left[\dfrac{\pi}{6},\dfrac{7\pi}{6}\right]\). 因此最大值在 \(x=0\) 处取得，为 \(f(0)=3\)；最小值在 \(x=\dfrac{5\pi}{6}\) 处取得，为 \(f\left(\dfrac{5\pi}{6}\right)=-2\sqrt3\).
\end{enumerate}
\end{solution}


\begin{problem}
如图，在平面直角坐标系 \(xOy\) 中，椭圆 \(E:\dfrac{x^2}{a^2}+\dfrac{y^2}{b^2}=1\)（\(a>b>0\)）的左，右焦点分别为 \(F_1\)，\(F_2\)，离心率为 \(\dfrac12\)，两准线之间的距离为 \(8\). 点 \(P\) 在椭圆 \(E\) 上，且位于第一象限，过点 \(F_1\) 作直线 \(PF_1\) 的垂线 \(l_1\)，过点 \(F_2\) 作直线 \(PF_2\) 的垂线 \(l_2\).

\begin{enumerate}
\item 求椭圆 \(E\) 的标准方程；
\item 若直线 \(l_1\)，\(l_2\) 的交点 \(Q\) 在椭圆 \(E\) 上，求点 \(P\) 的坐标.
\end{enumerate}

\FigureLayoutDeclare{img/2017/jiangsu/q17_fig1.png}{6.5cm}{333bp 105bp 123bp 0bp}
\bitmapfigure[max width=6.0cm]{img/2017/jiangsu/q17_fig1.png}
\end{problem}

\begin{solution}
\begin{enumerate}
\item 设椭圆的半焦距为 \(c\). 由离心率 \(\dfrac ca=\dfrac12\)，得右准线方程为 \(x=\dfrac{a^2}{c}=2a\)，所以两准线之间的距离为 \(4a=8\)，从而 \(a=2\)，\(c=1\). 因此 \(b^2=a^2-c^2=3\)，椭圆的标准方程为
\[
\dfrac{x^2}{4}+\dfrac{y^2}{3}=1.
\]
\item 由 \(a=2\)，\(c=1\)，得 \(F_1=(-1,0)\)，\(F_2=(1,0)\). 设 \(P=(x,y)\)，则 \(x>0\)，\(y>0\)，且 \(\dfrac{x^2}{4}+\dfrac{y^2}{3}=1\). 设 \(Q=(u,v)\). 因为 \(l_1\perp PF_1\)，\(l_2\perp PF_2\)，所以
\[
(u+1)(x+1)+vy=0,\qquad (u-1)(x-1)+vy=0.
\]
两式相减得 \(u=-x\)，代入第一式得 \(v=\dfrac{x^2-1}{y}\). 又因为 \(Q\) 在椭圆上，所以
\[
\dfrac{x^2}{4}+\dfrac{(1-x^2)^2}{3y^2}=1.
\]
由 \(y^2=3\left(1-\dfrac{x^2}{4}\right)\)，化简得
\[
\dfrac{x^2}{4}+\dfrac{4(1-x^2)^2}{9(4-x^2)}=1,\qquad 7x^4+40x^2-128=0.
\]
令 \(t=x^2\)，则 \(7t^2+40t-128=0\)，解得 \(t=\dfrac{16}{7}\) 或 \(t=-8\). 因为 \(x>0\)，所以 \(x=\dfrac{4\sqrt7}{7}\)，再由椭圆方程得 \(y=\dfrac{3\sqrt7}{7}\). 故 \(P\) 的坐标为 \(\left(\dfrac{4\sqrt7}{7},\dfrac{3\sqrt7}{7}\right)\).
\end{enumerate}
\end{solution}


\begin{problem}
如图，水平放置的正四棱柱形玻璃容器 \(I\) 和正四棱台形玻璃容器 \(II\) 的高均为 \(32\text{ cm}\)，容器 \(I\) 的底面对角线 \(AC\) 的长为 \(10\sqrt7\text{ cm}\)，容器 \(II\) 的两底面对角线 \(EG\)，\(E_1G_1\) 的长分别为 \(14\text{ cm}\) 和 \(62\text{ cm}\). 分别在容器 \(I\) 和容器 \(II\) 中注入水，水深均为 \(12\text{ cm}\). 现有一根玻璃棒 \(l\)，其长度为 \(40\text{ cm}\). （容器厚度，玻璃棒粗细均忽略不计）

\begin{enumerate}
\item 将 \(l\) 放在容器 \(I\) 中，\(l\) 的一端置于点 \(A\) 处，另一端置于侧棱 \(CC_1\) 上，求 \(l\) 没入水中部分的长度；
\item 将 \(l\) 放在容器 \(II\) 中，\(l\) 的一端置于点 \(E\) 处，另一端置于侧棱 \(GG_1\) 上，求 \(l\) 没入水中部分的长度.
\end{enumerate}

\bitmapinlinefigure[11.2cm][]{img/2017/jiangsu/q18_fig1.png}
\end{problem}

\begin{solution}
\begin{enumerate}
\item 设玻璃棒的另一端在侧棱 \(CC_1\) 上，且该端到容器 \(I\) 底面的距离为 \(h\). 由于容器 \(I\) 为直棱柱，玻璃棒的水平投影长度为底面对角线 \(AC=10\sqrt7\text{ cm}\)，所以
\[
(10\sqrt7)^2+h^2=40^2,
\]
解得 \(h=30\text{ cm}\). 玻璃棒从底面上的点 \(A\) 出发，水面高度为 \(12\text{ cm}\)，因此没入水中的长度为
\[
40\cdot\dfrac{12}{30}=16\text{ cm}.
\]
\item 设玻璃棒的另一端在侧棱 \(GG_1\) 上，设该端的高度为 \(32t\text{ cm}\)，其中 \(0\le t\le1\). 正四棱台的底面和顶面中心重合，且对应顶点的连线为侧棱，所以从底面顶点 \(G\) 到该端点的水平位移沿对角线方向. 底面对角线的一半为 \(7\text{ cm}\)，顶面对角线的一半为 \(31\text{ cm}\)，故点 \(E\) 到该端点的水平距离为
\[
7+(7+24t)=14+24t\text{ cm}.
\]
由玻璃棒长度为 \(40\text{ cm}\)，得
\[
(14+24t)^2+(32t)^2=40^2,
\]
即 \(400t^2+168t-351=0\). 因为 \(0\le t\le1\)，解得 \(t=\dfrac34\)，所以玻璃棒另一端的高度为 \(24\text{ cm}\). 水面高度为 \(12\text{ cm}\)，玻璃棒为直线段，长度与竖直方向的高度成正比，故没入水中的长度为
\[
40\cdot\dfrac{12}{24}=20\text{ cm}.
\]
\end{enumerate}
\end{solution}


\begin{problem}
对于给定的正整数 \(k\)，若数列 \(\{a_n\}\) 满足：
\(a_{n-k}+a_{n-k+1}+\allowbreak\cdots+\allowbreak a_{n-1}\)\allowbreak
\(+a_{n+1}+\allowbreak\cdots+\allowbreak a_{n+k-1}+a_{n+k}\)\allowbreak
\(=2ka_n\)，
对任意正整数 \(n\)（\(n>k\)）总成立，则称数列 \(\{a_n\}\) 是“\(P(k)\)数列”.

\begin{enumerate}
\item 证明：等差数列 \(\{a_n\}\) 是“\(P(3)\)数列”；
\item 若数列 \(\{a_n\}\) 既是“\(P(2)\)数列”，又是“\(P(3)\)数列”，证明：\(\{a_n\}\) 是等差数列.
\end{enumerate}
\end{problem}

\begin{solution}
\begin{enumerate}
\item 设等差数列 \(\{a_n\}\) 的公差为 \(d\)，则对 \(j=1,2,3\)，都有 \(a_{n-j}+a_{n+j}=2a_n\). 因此
\[
a_{n-3}+a_{n-2}+a_{n-1}+a_{n+1}+a_{n+2}+a_{n+3}=6a_n,
\]
所以 \(\{a_n\}\) 是 \(P(3)\) 数列.
\item 对任意 \(n>3\)，将 \(P(3)\) 的等式减去 \(P(2)\) 的等式，得
\[
a_{n-3}+a_{n+3}=2a_n.
\]
令 \(u=a_4-a_1\)，\(v=a_5-a_2\)，\(w=a_6-a_3\). 由上式递推可得，对任意非负整数 \(q\)，
\[
a_{3q+1}=a_1+qu,\qquad a_{3q+2}=a_2+qv,\qquad a_{3q+3}=a_3+qw.
\]
在 \(P(2)\) 等式中取 \(n=3q+3\)，有
\[
(a_1+qu)+(a_2+qv)+(a_1+(q+1)u)+(a_2+(q+1)v)=4(a_3+qw),
\]
比较一次项系数得 \(u+v=2w\). 取 \(n=3q+1\)（\(q\ge1\)），有
\[
(a_2+(q-1)v)+(a_3+(q-1)w)+(a_2+qv)+(a_3+qw)=4(a_1+qu),
\]
比较一次项系数得 \(v+w=2u\). 取 \(n=3q+2\)（\(q\ge1\)），有
\[
(a_3+(q-1)w)+(a_1+qu)+(a_3+qw)+(a_1+(q+1)u)=4(a_2+qv),
\]
比较一次项系数得 \(w+u=2v\). 因而
\[
u+v=2w,\qquad v+w=2u,\qquad w+u=2v.
\]
由前两式相减得 \(u=w\)，再由任一式得 \(u=v=w\). 记其公共值为 \(\delta\). 在 \(P(2)\) 等式中取 \(n=5\)，得
\[
a_3+a_4+a_6+a_7=4a_5.
\]
代入 \(a_4=a_1+\delta\)，\(a_5=a_2+\delta\)，\(a_6=a_3+\delta\)，\(a_7=a_1+2\delta\)，可得 \(a_1+a_3=2a_2\). 令 \(d=a_2-a_1\)，则 \(a_3=a_1+2d\). 再在 \(P(2)\) 等式中取 \(n=3\)，得
\[
a_1+a_2+a_4+a_5=4a_3,
\]
代入上述各式可得 \(\delta=3d\). 因而对三个剩余类分别有
\[
a_{3q+1}=a_1+3qd,\qquad a_{3q+2}=a_1+(3q+1)d,\qquad a_{3q+3}=a_1+(3q+2)d.
\]
这正说明 \(a_n=a_1+(n-1)d\)，所以 \(\{a_n\}\) 是等差数列.
\end{enumerate}
\end{solution}


\begin{problem}
已知函数 \(f(x)=x^3+ax^2+bx+1\)（\(a>0\)，\(b\in\R\)）有极值，且导函数 \(f'(x)\) 的极值点是 \(f(x)\) 的零点. （极值点是指函数取极值时对应的自变量的值）

\begin{enumerate}
\item 求 \(b\) 关于 \(a\) 的函数关系式，并写出定义域；
\item 证明：\(b^2>3a\)；
\item 若 \(f(x)\)，\(f'(x)\) 这两个函数的所有极值之和不小于 \(-\dfrac72\)，求 \(a\) 的取值范围.
\end{enumerate}
\end{problem}

\begin{solution}
\begin{enumerate}
\item \(f'(x)=3x^2+2ax+b\) 的极值点为 \(x=-\dfrac a3\). 由该点是 \(f(x)\) 的零点，得
\[
0=f\left(-\dfrac a3\right)=-\dfrac{a^3}{27}+\dfrac{a^3}{9}-\dfrac{ab}{3}+1=\dfrac{2a^3}{27}-\dfrac{ab}{3}+1.
\]
因为 \(a>0\)，所以
\[
b=\dfrac{2a^3+27}{9a}=\dfrac{2a^2}{9}+\dfrac3a.
\]
又因为三次函数 \(f(x)\) 有极值，所以二次方程 \(f'(x)=0\) 有两个不相等的实根，即 \(a^2-3b>0\). 代入上式得 \(a^2>\dfrac{2a^2}{3}+\dfrac9a\)，从而 \(a^3>27\). 因此定义域为 \(a>3\).
\item 由 \(a>3\)，有
\[
b^2-3a=\left(\dfrac{2a^2}{9}+\dfrac3a\right)^2-3a=\dfrac{(a^3-27)(4a^3-27)}{81a^2}>0.
\]
故 \(b^2>3a\).
\item 由
\[
f'(x)=3\left(x+\dfrac a3\right)^2+b-\dfrac{a^2}{3},
\]
知 \(f'(x)\) 的极小值为 \(b-\dfrac{a^2}{3}\). 设 \(f'(x)=0\) 的两个根为 \(r_1\)，\(r_2\)，则它们是 \(f(x)\) 的两个极值点，且
\[
r_1+r_2=-\dfrac{2a}{3},\qquad r_1r_2=\dfrac b3.
\]
于是
\[
\begin{aligned}
f(r_1)+f(r_2)&=(r_1+r_2)^3-3r_1r_2(r_1+r_2)+a\left((r_1+r_2)^2-2r_1r_2\right)+b(r_1+r_2)+2\\
&=\dfrac{4a^3}{27}-\dfrac{2ab}{3}+2=2f\left(-\dfrac a3\right)=0.
\end{aligned}
\]
所以两个函数的所有极值之和为 \(b-\dfrac{a^2}{3}=\dfrac{27-a^3}{9a}\). 由题意
\[
\dfrac{27-a^3}{9a}\ge-\dfrac72,
\]
即 \(2a^3-63a-54\le0\). 因为
\[
2a^3-63a-54=(a-6)(2a^2+12a+9),
\]
且 \(a>3\) 时 \(2a^2+12a+9>0\)，所以 \(3<a\le6\). 故 \(a\) 的取值范围是 \((3,6]\).
\end{enumerate}
\end{solution}

\section{附加题}

\begin{problem}[A]
如图，\(AB\) 为半圆 \(O\) 的直径，直线 \(PC\) 切半圆 \(O\) 于点 \(C\)，\(AP\perp PC\)，\(P\) 为垂足. 求证：

\begin{enumerate}
\item \(\angle PAC=\angle CAB\)；
\item \(AC^2=AP\cdot AB\).
\end{enumerate}

\bitmapfigure[max width=6.0cm]{img/2017/jiangsu/q21_fig1.png}
\end{problem}

\begin{solution}
\begin{enumerate}
\item 因为 \(PC\) 是半圆在点 \(C\) 处的切线，所以 \(OC\perp PC\). 又已知 \(AP\perp PC\)，故 \(AP\parallel OC\). 在等腰三角形 \(AOC\) 中，\(OA=OC\)，所以 \(\angle OAC=\angle OCA\). 由于 \(A\)，\(O\)，\(B\) 共线，且射线 \(AO\) 与射线 \(AB\) 同向，故 \(\angle CAB=\angle OAC\). 又由 \(AP\parallel OC\)，得 \(\angle PAC=\angle OCA\). 因此 \(\angle PAC=\angle CAB\).
\item 因为 \(AB\) 是半圆的直径，所以 \(\angle ACB=90^\circ\). 又 \(AP\perp PC\)，所以 \(\angle APC=90^\circ\). 结合第（1）问，得 \(\triangle APC\sim\triangle ACB\). 因而
\[
\dfrac{AP}{AC}=\dfrac{AC}{AB},
\]
即 \(AC^2=AP\cdot AB\).
\end{enumerate}
\end{solution}


\reuseproblemnumber
\begin{problem}[B]
已知矩阵 \(A=\begin{pmatrix}0&1\\1&0\end{pmatrix}\)，\(B=\begin{pmatrix}1&0\\0&2\end{pmatrix}\).

\begin{enumerate}
\item 求 \(AB\)；
\item 若曲线 \(C_1:x^2+4y^2=8\) 在矩阵 \(AB\) 对应的变换作用下得到另一曲线 \(C_2\)，求 \(C_2\) 的方程.
\end{enumerate}
\end{problem}

\begin{solution}
\begin{enumerate}
\item
\[
AB=\begin{pmatrix}0&1\\1&0\end{pmatrix}\begin{pmatrix}1&0\\0&2\end{pmatrix}=\begin{pmatrix}0&2\\1&0\end{pmatrix}.
\]
\item 设曲线 \(C_1\) 上任意一点 \((x,y)\) 在矩阵 \(AB\) 对应的变换下变为 \((X,Y)\)，则
\[
\begin{pmatrix}X\\Y\end{pmatrix}=AB\begin{pmatrix}x\\y\end{pmatrix}=\begin{pmatrix}2y\\x\end{pmatrix}.
\]
所以 \(x=Y\)，\(y=\dfrac X2\). 代入 \(C_1\) 的方程得 \(Y^2+4\left(\dfrac X2\right)^2=8\)，即 \(X^2+Y^2=8\). 因此将坐标字母改回 \(x\)，\(y\)，得 \(C_2:x^2+y^2=8\).
\end{enumerate}
\end{solution}


\reuseproblemnumber
\begin{problem}[C]
在平面直角坐标系 \(xOy\) 中，已知直线 \(l\) 的参数方程为
\(\begin{cases}
x=-8+t,\\
y=\dfrac{t}{2}
\end{cases}\)（\(t\) 为参数），曲线 \(C\) 的参数方程为
\(\begin{cases}
x=2s^2,\\
y=2\sqrt2\,s
\end{cases}\)（\(s\) 为参数）. 设 \(P\) 为曲线 \(C\) 上的动点，求点 \(P\) 到直线 \(l\) 的距离的最小值.
\end{problem}

\begin{solution}
直线 \(l\) 的参数方程消去参数 \(t\) 得 \(x-2y+8=0\). 设曲线 \(C\) 上的点为 \(P=(2s^2,2\sqrt2s)\)，则点 \(P\) 到直线 \(l\) 的距离为
\[
d(s)=\dfrac{\left|2s^2-4\sqrt2s+8\right|}{\sqrt5}=\dfrac{2\left((s-\sqrt2)^2+2\right)}{\sqrt5}.
\]
当且仅当 \(s=\sqrt2\) 时，距离取得最小值 \(\dfrac4{\sqrt5}\). 此时 \(P=(4,4)\)，所以所求最小值为 \(\dfrac4{\sqrt5}\).
\end{solution}


\reuseproblemnumber
\begin{problem}[D]
已知 \(a\)，\(b\)，\(c\)，\(d\) 为实数，且 \(a^2+b^2=4\)，\(c^2+d^2=16\)，证明 \(ac+bd\le 8\).
\end{problem}

\begin{solution}
由柯西不等式，
\[
ac+bd\le\sqrt{a^2+b^2}\sqrt{c^2+d^2}=\sqrt4\sqrt{16}=8.
\]
所以 \(ac+bd\le8\).
\end{solution}




\begin{problem}
如图，在平行六面体 \(ABCD-A_1B_1C_1D_1\) 中，\(AA_1\perp\) 平面 \(ABCD\)，且 \(AB=AD=2\)，\(AA_1=\sqrt3\)，\(\angle BAD=120^\circ\).

\begin{enumerate}
\item 求异面直线 \(A_1B\) 与 \(AC_1\) 所成角的余弦值；
\item 求二面角 \(B-A_1D-A\) 的正弦值.
\end{enumerate}

\bitmapfigure[max width=4.6cm]{img/2017/jiangsu/q25_fig1.png}
\end{problem}

\begin{solution}
建立空间直角坐标系：原点取为 \(A\)，\(x\) 轴沿 \(AB\)，\(y\) 轴在平面 \(ABCD\) 内且垂直于 \(AB\)，\(z\) 轴沿 \(AA_1\). 由已知条件，
\[
A=(0,0,0),\quad B=(2,0,0),\quad D=(-1,\sqrt3,0),\quad C=(1,\sqrt3,0),
\]
\[
A_1=(0,0,\sqrt3),\quad C_1=(1,\sqrt3,\sqrt3).
\]
\begin{enumerate}
\item 直线 \(A_1B\) 和 \(AC_1\) 的方向向量可分别取为 \((2,0,-\sqrt3)\) 和 \((1,\sqrt3,\sqrt3)\). 因此它们的数量积为 \(-1\)，长度均为 \(\sqrt7\). 异面直线所成角取锐角，故其余弦值为
\[
\dfrac{|-1|}{\sqrt7\sqrt7}=\dfrac17.
\]
\item 平面 \(BA_1D\) 的一个法向量为
\[
\bs{n}_1=(2,0,-\sqrt3)\times(-1,\sqrt3,-\sqrt3)=(3,3\sqrt3,2\sqrt3),
\]
平面 \(AA_1D\) 的一个法向量为
\[
\bs{n}_2=(0,0,-\sqrt3)\times(-1,\sqrt3,-\sqrt3)=(3,\sqrt3,0).
\]
设二面角 \(B-A_1D-A\) 的大小为 \(\theta\)，则
\[
\left|\cos\theta\right|=\dfrac{|\bs{n}_1\cdot\bs{n}_2|}{|\bs{n}_1||\bs{n}_2|}=\dfrac{18}{4\sqrt3\cdot2\sqrt3}=\dfrac34.
\]
所以 \(\sin\theta=\sqrt{1-\dfrac9{16}}=\dfrac{\sqrt7}{4}\).
\end{enumerate}
\end{solution}


\begin{problem}
已知一个口袋有 \(m\) 个白球，\(n\) 个黑球（\(m,n\in \N^*\)，\(n\ge 2\)），这些球除颜色外全部相同. 现将口袋中的球随机地逐个取出，并放入如图所示的编号为 \(1\)，\(2\)，\(3\)，\(\ldots\)，\(m+n\) 的抽屉内，其中第 \(k\) 次取出的球放入编号为 \(k\) 的抽屉（\(k=1\)，\(2\)，\(3\)，\(\ldots\)，\(m+n\)）.

\begin{center}
\begin{tabular}{|c|c|c|c|c|}
\hline
\(1\) & \(2\) & \(3\) & \(\ldots\) & \(m+n\) \\
\hline
\end{tabular}
\end{center}

\begin{enumerate}
\item 试求编号为 \(2\) 的抽屉内放的是黑球的概率 \(p\)；
\item 随机变量 \(X\) 表示最后一个取出的黑球所在抽屉编号的倒数，\(E(X)\) 是 \(X\) 的数学期望，证明 \(E(X)<\dfrac{n}{(m+n)(n-1)}\).
\end{enumerate}
\end{problem}

\begin{solution}
\begin{enumerate}
\item 在 \(m+n\) 个抽屉中，\(n\) 个黑球的位置等可能地从所有位置组合中产生，所以编号为 \(2\) 的抽屉内放黑球的概率为
\[
p=\dfrac{n}{m+n}.
\]
\item 令 \(N=m+n\)，设最后一个黑球所在的抽屉编号为 \(K\)，则 \(X=\dfrac1K\). 当 \(k=n,n+1,\ldots,N\) 时，事件 \(K=k\) 等价于第 \(k\) 个抽屉放黑球，且前 \(k-1\) 个抽屉中有其余的 \(n-1\) 个黑球，因此
\[
P(K=k)=\dfrac{\mathrm C_{k-1}^{n-1}}{\mathrm C_N^n}.
\]
于是
\[
E(X)=\dfrac1{\mathrm C_N^n}\sum_{k=n}^{N}\dfrac{\mathrm C_{k-1}^{n-1}}{k}.
\]
由于 \(n\ge2\)，对每个 \(k\ge n\) 都有 \(\dfrac1k<\dfrac1{k-1}\)，从而
\[
E(X)<\dfrac1{\mathrm C_N^n}\sum_{k=n}^{N}\dfrac{\mathrm C_{k-1}^{n-1}}{k-1}=\dfrac1{(n-1)\mathrm C_N^n}\sum_{k=n}^{N}\mathrm C_{k-2}^{n-2}.
\]
由组合数的累加公式，
\[
\sum_{k=n}^{N}\mathrm C_{k-2}^{n-2}=\mathrm C_{N-1}^{n-1}.
\]
因此
\[
E(X)<\dfrac{\mathrm C_{N-1}^{n-1}}{(n-1)\mathrm C_N^n}=\dfrac{n}{N(n-1)}=\dfrac{n}{(m+n)(n-1)}.
\]
\end{enumerate}
\end{solution}
